\chapter*{\chapterstyle{Géométrie {:} Espaces affines}}
Soit \(\mathscr{E}\) un ensemble non-vide et \(V\) un \(\R\)-espace vectoriel de dimension \(n\) finie.\+
On appelle les éléments de \(\mathscr{E}\) des \textbf{points}.

\subsection*{\subsecstyle{Définition {:}}}
On appelle \textbf{espace affine de direction} \(V\) le couple \((\mathscr{E}, +)\) avec:
\[
   \begin{aligned}
      + : &&\mathscr{E} \times V &\longrightarrow \mathscr{E}\\
      &&(\,A, \, u\,) &\longmapsto A + u
   \end{aligned}
\]
La loi + doit vérifer les axiomes suivants 
\footnote{Une action d'un groupe \((G, \star)\) sur \(E\) est une application \(+\) de \(G \times E \) dans \(E\) qui vérifie: \[
   \forall g, g' \in G  \; , \; \forall x \in E \; ; \; x + (g \star g') = (x + g) \star g'
\]}:
\multiparbox{-0.8em}{0.8}{
    \begin{align*}
      &\textbf{Existence d'un neutre} && \forall A \in \mathscr{E} \; ; \; A + 0_V = A\\
      &\textbf{Action de groupe} && \forall A \in \mathscr{E} \; , \; \forall u,v \in V \; ; \; A + (u + v) = (A + u) + v \\
      &\textbf{Unicité du translaté} && \forall A, B \in \mathscr{E} \; , \; \exists !u \in V \; ; \; A + u = B 
    \end{align*}
}
Etant donné deux points \(A, B \in \mathscr{E}\), on note alors \(\overrightarrow{AB}\) l'unique vecteur \(u\) qui vérifie:
\[
   \mathbox{A + u = B}   
\]

\subsection*{\subsecstyle{Propriétés {:}}}
Si \(\mathscr{E}\) est un espace affine de direction \(V\), on a alors les propriétés suivantes:
\multiparbox{-0.3em}{0.4}{
    \begin{align*}
      &\textbf{Relation de Chasles} && \overrightarrow{AB} + \overrightarrow{BC} = \overrightarrow{AC} \\
      &\textbf{Existence d'un neutre} && \overrightarrow{AB} + \overrightarrow{BA} = O_V
    \end{align*}
}

\subsection*{\subsecstyle{Repères {:}}}
Un repère de \(\mathscr{E}\) est un couple \(\mathscr{R} = (O, \mathscr{B})\) formé d'un point \(O \in \mathscr{E}\) et d'une base \(\mathscr{B} = (e_1, \ldots, e_n)\) de \(V\). On appelle alors le point \(O\) \textbf{origine} du repère, et les vecteurs \(e_1, \ldots, e_n\) \textbf{vecteurs de base} du repère.\<

Soit \(i\in\inticc{1}{n}\), alors pour tout point \(A \in \mathscr{E}\), on appelle \textbf{coordonées} de \(A\) dans le repère \(\mathscr{R}\), les composantes \((x_i)_i\) du vecteur qui représente la translation de \(O\) vers \(A\) dans la base \(\mathscr{B}\), et on a la caractérisation élémentaire suivante:
\[
   \overrightarrow{OA} = x_1e_1 + \ldots + x_ne_n 
\]

\subsection*{\subsecstyle{Cas particulier des espaces vectoriels {:}}}
Il est important de noter que le couple \((V, +)\) forme un espace affine de \textbf{direction lui-même}. En effet, si on considère un élément de \(V\) comme un point, alors le couple  \((V, +)\) forme un espace affine de direction \(V\) et la loi \(+\) est alors exactement la loi de composition interne de \(V\).\<

L'unique vecteur \(u\) tel que \(A + u = B\) est exactement \(B - A\) (ici \(A, B\) sont des points de \(V\) qui se trouvent être des vecteurs dans ce cas particulier). \<\<

\chapter*{\chapterstyle{Géométrie {:} Espaces euclidiens}}
Soit \(E\) un \(\R\)-espace vectoriel de dimension \(n\) finie. On considère \(E\) avec sa structure d'espace affine canonique.
On appelle \textbf{espace euclidien} un tel espace vectoriel si il est muni d'un \textbf{produit scalaire}.\<

Soit \(u, v\) deux vecteurs de \(E\) et \((u_i)_i\) pour \(i \in \inticc{1}{n}\) les composantes de \(u\) dans la base canonique de \(E\).

\subsection*{\subsecstyle{Vecteurs {:}}}
Soit deux points \(A, B\), on note \(\overrightarrow{AB}\) l'unique vecteur \(u\) de \(E\) tel que \(A + u = B\), on appelle alors \(A\) le \textbf{point initial} et \(B\) le \textbf{point final}. \<

On peut alors considérer l'application \(\phi\) qui associe aux points \(A, B\) le vecteur \(\overrightarrow{AB}\) en \textbf{position standard}, et on a:
\[
   \begin{aligned}
      \phi: E \times E &\longrightarrow E\\
      (A, B) &\longmapsto B - A
   \end{aligned}
\]
Ici, \(A, B\) sont des points qui trouvent aussi être des vecteurs, car l'espace affine sous-jacent est un espace vectoriel, et donc la soustractrion ci-dessus est bien définie.


\subsection*{\subsecstyle{Produit scalaire {:}}}
Soit \(\phi : E^2 \longrightarrow \R\) une forme \textbf{bilinéaire}. On dit que l'application \(\phi\) est un \textbf{produit scalaire} si et seulement si il vérifie:

\multiparbox{-0.3em}{0.4}{
    \begin{align*}
      &\textbf{Symétrique} &&\dotproduct{u}{v} = \dotproduct{v}{u} \\
      &\textbf{Définie} &&\dotproduct{u}{u} = 0 \implies u = 0 \\
      &\textbf{Positive} &&\dotproduct{u}{v} \geq 0
    \end{align*}
}

Par exemple, si on considère \(E = \R^n\), alors on peut définir le produit scalaire canonique de \(\R^n\) comme ci-dessous:
\[
   \phi : (u, v) \longmapsto \sum_{k=1}^{n} u_iv_i 
\]
Par ailleurs, on a \textbf{l'inégalité de Cauchy-Schwartz}:
\[
    |\dotproduct{u}{v}| \leq \vectNorm{u} \vectNorm{v}
\]

\subsection*{\subsecstyle{Norme{:}}}
On appelle \textbf{norme} une application \(\vectNorm{\cdot}\) de E vers \(\R\) qui vérifie:
\multiparbox{-0.3em}{0.4}{
    \begin{align*}
      &\textbf{Séparation} &&\vectNorm{u} = 0 \implies u = 0\\
      &\textbf{Absolue Homogénéité} &&\vectNorm{\lambda u} = |\lambda| \vectNorm{u}\\
      &\textbf{Sous-Additivité} &&\vectNorm{u + v} \leq \vectNorm{u} + \vectNorm{v}
    \end{align*}
}
Un espace vectoriel muni d'une norme est alors appelé \textbf{espace vectoriel normé}. Dans le cas d'un espace euclidien, on définit la norme euclidienne usuelle:
\[
   \vectNorm{u} = \sqrt{\sum_{i=1}^n u_i^2}
\]
Cette norme induit alors une \textbf{métrique} sur \(E\)  comme suit:
\[
   \begin{aligned}
      d: E \times E &\longrightarrow E\\
      (u, v) &\longmapsto \vectNorm{v - u}
   \end{aligned}
\]
Muni de cette métrique, l'espace euclidien \(E\) est alors un \textbf{espace métrique homogène}.

\subsection*{\subsecstyle{Angles{:}}}
Dans un espace euclidien, on définit \textbf{l'angle} \(\theta := \angle AOC\) comme étant l'angle entre les bipoints \(OA\) et \(OC\) et on a alors la formule de \textbf{Pythagore généralisée}:
\[
   \vectNorm{BC}^2 = \vectNorm{AB}^2 + \vectNorm{AC}^2 - 2\vectNorm{AB}\vectNorm{AC}\cos(\theta)
\]
Puis par suite on peut montrer la caractérisation du produit scalaire:
\[
    \mathbox{\dotproduct{u}{v} = \vectNorm{u} \vectNorm{v} \cos{(\theta)}}
\]

\subsection*{\subsecstyle{Composantes scalaires et projections{:}}}
Soit \(\alpha\) l'angle entre les deux vecteurs \(u\) et \(v\).
On appelle \textbf{composante scalaire} de \(u\) par rapport à \(v\) le scalaire:
\[
   \text{comp}_v(u) = \vectNorm{u}\cos(\alpha)   
\]
On définit alors la \textbf{projection orthogonale} de \(u\) sur \(v\):
\[
   \mathbox{\text{proj}_v(u) = \text{comp}_v(u)\frac{v}{\vectNorm{v}}}   
\]
Visuellement, si on considère deux vecteurs de \(E = \R^2\), on a:
\begin{center}
   \begin{tikzpicture}[
      >=stealth,scale=1,line cap=round,
      bullet/.style={circle,inner sep=0.5pt,fill}
   ]
      \draw[->] (-1,0) -- (9,0);
      \draw[->] (0,-1) -- (0, 5.5);
   
      \draw foreach \X in {3,6}{
         (\X, 0.1) -- (\X, -0.1) node[below]{$\X$}
      };
      \draw foreach \Y in {2,4}{
         (0.1,\Y) -- (-0.1,\Y) node[left]{$\Y$}
      };
   
      \draw[->]  (0,0) coordinate (O) -- (3,4) coordinate (B) node[above]{$u$};
      \draw[->]  (0,0) coordinate (O) -- (6,2) coordinate (A) node[above]{$v$};

      \draw[dashed, color=DarkBlue1] (3, 4) -- (3.9, 1.3);

      \draw[->, line width=2pt, color=DarkBlue1]  (0,0) -- (3.9, 1.3);
      \path (2,1.2) node[color=DarkBlue1]{$\text{proj}_v(u)$};
      \draw[solid, line width=2pt, color=BrightBlue1]  (0,0) -- (4.111, 0) node[above]{$\text{comp}_v(u)$};
      \draw pic["$\alpha$", draw, gray, angle radius=1cm] {angle = A--O--B};

   \end{tikzpicture}   
\end{center}
On remarque alors que \textbf{la composante scalaire de \(u\) est exactement la norme du projeté de \(u\) sur \(v\)}

\subsection*{\subsecstyle{Orientation{:}}}
On se place dans \(E = \R^3\), soit deux vecteur \(u, v\) non parallèles, on appelle \textbf{orientation directe} l'orientation de l'espace suivante \textbf{la règle de la main droite}, et orientation indirecte celle suivant la règle de la main gauche.
\pagebreak

\subsection*{\subsecstyle{Produit vectoriel{:}}}
On définit le \textbf{produit vectoriel} de deux vecteurs \(u, v\) comme l'unique vecteur \(w\) qui vérifie:
\multiparbox{-0.3em}{0.4}{
    \begin{center}
      $\vectNorm{w} = \vectNorm{u}\vectNorm{v}|\sin(\theta)|$\\
      $w$ \textbf{est orthogonal à }$u$\textbf{ et à }$v$\\
      \textbf{La base } $(u, v, w)$ \textbf{ est directe}
    \end{center}
}

C'est une application de \(E \times E \longrightarrow E\) qui vérfie:
\multiparbox{-0.5em}{0.6}{
    \begin{align*}
      &\textbf{Homogénéité} && \lambda_1 u \times v = \lambda_1 (u \times v)\\
      &\textbf{Anticommutativité} && u \times v = - (v \times u)\\
      &\textbf{Distributivité} && u \times (v + w) = u \times v + u \times w
    \end{align*}
}
Néanmoins, le produit vectoriel n'est en général \textbf{pas associatif}. En particulier, on a:
\[
    u \times (v \times w) = \dotproduct{u}{w}v - \dotproduct{u}{v}w
\]

\underline{Calcul d'un produit vectoriel:}\<

On note \(\big|M\big|\)  le déterminant d'une matrice, alors on a une méthode efficace pour calculer des produits vectoriels:
\[
   \begin{cases}
      u := (u_1, u_2, u_3)\\ 
      v := (v_1, v_2, v_3)
   \end{cases} \implies 
   u \times v = \Biggl(\mdetTwo{u_2}{u_3}{v_2}{v_3} ; -\mdetTwo{u_1}{u_3}{v_1}{v_3} ; \mdetTwo{u_1}{u_2}{v_1}{v_2}\Biggl)
\]
\subsection*{\subsecstyle{Volumes \& Produit triple{:}}}
On peut remarque que la norme du vecteur resultant d'un produit vectoriel est \textbf{l'aire du parallélogramme} formé par \(u\) et \(v\).\<

On peut alors montrer que le produit \(\dotproduct{u \times v}{w}\) est \textbf{le volume du parallélépipède} formé par \(u, v, w\). En effet, on a:
\[
   \dotproduct{u \times v}{w} = \vectNorm{u \times v} \vectNorm{w} \cos(\alpha) = \vectNorm{u \times v}h
\]
Avec \(\vectNorm{u \times v}\) qui est l'aire de la base est \(h\) la hauteur du parallélépipède.

\begin{center}
   \begin{tikzpicture}[
      >=stealth,scale=1,line cap=round,
      bullet/.style={circle,inner sep=0.5pt,fill}
   ]
      \coordinate (A) at (0, 0, 0);
      \coordinate (B) at (8, 0, 0);
      \coordinate (C) at (8, 0, -8);
      \coordinate (D) at (0, 0, -8);
      \coordinate (E) at (0, 2, -10);
      \coordinate (F) at (0, 2, -2);
      \coordinate (G) at (8, 2, -2);
      \coordinate (H) at (8, 2, -10);

      \draw[->] (A) -- (B) node[below right]{\(u\)};
      \draw[->] (A) -- (D) node[above left]{\(v\)};
      \draw[->] (A) -- (F) node[below right]{\(w\)};

      \draw[dashed, color=gray] (B) -- (C) -- (D) -- (E) -- (F) -- (G) -- (H) -- (E);
      \draw[dashed, color=gray] (B) -- (G);
      \draw[dashed, color=gray] (C) -- (H);
      \draw[dashed, color=black!10, fill=BrightBlue1, opacity=0.2] (E) -- (F) -- (G) -- (H) -- cycle;
      \draw[dashed, color=black!10, fill=BrightBlue1, opacity=0.15] (A) -- (B) -- (C) -- (D) -- cycle;
      \draw[dashed, color=black!10, fill=BrightBlue1, opacity=0.2] (A) -- (B) -- (G) -- (F) -- cycle;
      \draw[dashed, color=black!10, fill=BrightBlue1, opacity=0.2] (B) -- (C) -- (H) -- (G) -- cycle;

      \draw[->] (A) -- (0, 4, 0) coordinate (X) node[below right]{\(u \times v\)};
      \draw[gray] pic ["$\alpha$", draw, angle radius=2cm] {angle = F--A--X};

      \draw[->, color = DarkBlue3, line width = 1.2pt] (0, 0.01, 0) -- (0, 2.7, 0) coordinate (X) node[left]{\(h\)};
   \end{tikzpicture}   
\end{center}
Un autre manière de déterminer le volume d'un parallélépipède est de calculer le déterminant \(\mdetThree{u_1}{u_2}{u_3}{v_1}{v_2}{v_3}{w_1}{w_2}{w_3}\)

\chapter*{\chapterstyle{Géométrie {:} Droites \& Plans}}
Dans l'espace euclidien \(\R^3\) , nous savons qu'on peut définir des droites vectorielles, des plans ou des volumes, dans ce chapitre nous allons examiner comment les exprimer mathématiquement, calculer des distances entre ces figures ou des intersections. \<

Dans la suite, on considère un point \(P_0\) de coordonées \((x_0, y_0, z_0)\) et une vecteur \(\vv{v}\) de composantes \((a, b, c)\).

\subsection*{\subsecstyle{Droites{:}}}
Une droite de l'espace en trois dimensions est caractérisée par la donnée d'un \textbf{point} sur la droite, et d'un \textbf{vecteur directeur} de cette droite.\<

Formellement, la droite \(\mathscr{D}\) passant par \(A\) de direction \(\vv{v}\) est l'ensemble des points \(P\) tels que le vecteur 
\(\vv{PQ}\) est paralléle à \(\vv{v}\).\+
Si on considère une fonction \(\vv{f}(t)\) qui donne la position d'un point de la droite \(\mathscr{D}\), avec \(\vv{f}(0)\) la position de \(P_0\), on peut alors montrer que les points de la droite vérifient l'équation:
\[
   \mathbox{f(t) = \vv{f}(0) + t\vv{v}}   
\]
Si on décompose les vecteurs obtenus avec cette equation en composantes, on obtient que \textbf{les coordonées} des points de la droite vérifient:
\[
   P = \begin{cases} 
      x = x_0 + at\\
      y = y_0 + bt\\
      z = z_0 + ct
   \end{cases} 
\]
La paramétrisation n'est évidemment pas unique.

\subsection*{\subsecstyle{Plans{:}}}
Un plan de l'espace en trois dimensions est caractérisée par la donnée d'un \textbf{point} sur le plan, et d'un \textbf{vecteur normal} au plan.\<

Formellement, le plan \(\mathscr{P}\) passant par \(A\) de vecteur normal \(\vv{v}\) est l'ensemble des points \(P\) tels que le vecteur 
\(\vv{AP}\) est othogonal à \(\vv{v}\).\+
On en déduit que les points du plan \(\mathscr{P}\) vérifient l'équation vectorielle:
\[
   \mathbox{
      \dotproduct{\vv{v}}{\vv{AP}} = 0  
   }
\]
Si on déploie le produit sclaire en composantes, on obtient que \textbf{les coordonées} des points du plan vérifient:
\[
   a(x - x_0) + b(y - y_0) + c(z - z_0) = 0 \Longleftrightarrow ax + by + cz = C
\]
On peut alors considérer le lien avec l'espace vectoriel sous-jacent:
\begin{center}
   \textit{
      Les plans et les droites de l'espace en trois dimensions sont exactement les translations des droites et plans linéaires de l'espace vectoriel sous-jacent. 
   }
\end{center}

\subsection*{\subsecstyle{Distances \& Angles{:}}}
Nous sommes amenés à calculer des distances et des angles en géométrie euclidienne, par exemple calculer la distance entre un point \(P\) et une droite \(\mathscr{D}\).
Si nous connaissons un point \(A\) de la droite et un vecteur directeur \(\vv{v}\) de la droite, il suffit alors de calculer:
\begin{flalign*}
   &&\mathbox{\vectNorm{\vv{AP}}\sin(\alpha)} &\shorteqnote{\(\bigstar\)}
\end{flalign*}


L'apparition du sinus s'explique par le fait que la distance entre \(P\) et \(\mathscr{D}\) est \textbf{invariante par rotation}, et alors on peut se ramener à un problème trigonométrique sur un cercle de rayon \(\vectNorm{\vv{AP}}\).\+
Visuellement, on a:
\begin{center}
   \begin{tikzpicture}[
      >=stealth,scale=1,line cap=round,
      bullet/.style={circle,inner sep=0.5pt,fill}
   ]
      \coordinate (O) at (0,0);
      \coordinate (A) at (6,2);
      \coordinate (A') at (4.5, 1.5);
      \coordinate (P) at (3,4);

      \coordinate (V') at (-1.5, 4.5);
      \coordinate (Proj1) at (3.9, 1.3);
      \coordinate (Proj2) at (-0.9, 2.7);

      \draw[]  (-2, -2/3)  -- (8, 8/3) coordinate (A) node[above]{$\mathscr{D}$};

      \draw[->, color=DarkBlue1, line width = 1pt]  (O)  node[below]{$A$} -- (P) node[above]{$P$};
      \draw[->, color=DarkBlue1, line width = 1pt]  (O) -- (A') node[above]{$v$};
      \draw[->, color=DarkBlue1, line width = 1pt]  (O) -- (V') node[above]{$v^\perp$};
      \draw[solid, line width=2pt, color=BrightBlue1] (P) -- (Proj1);

      \draw node[color=BrightBlue1] at (-1.8, 1.7) {$\vectNorm{\vv{AP}}\sin(\alpha)$};

      \draw[solid, line width=2pt, color=BrightBlue1]  (O) -- (Proj2);
      \draw[dashed, color=DarkBlue1] (Proj2) -- (P);

      \draw pic["$\alpha$", draw, gray, angle radius=1cm] {angle = A--O--P};
      \draw pic["$\alpha'$", draw, gray, angle radius=1.2cm] {angle = P--O--Proj2};

      \RightAngle{(-2, -2/3)}{(0,0)}{(-1.5, 4.5)};
      \tkzMarkSegment[color=black, pos=0.5, size=5pt, mark=||](O, Proj2)
      \tkzMarkSegment[color=black, pos=0.5, size=5pt,mark=||](P, Proj1)

   \end{tikzpicture}   
\end{center}
On remarque alors une symétrie conceptuelle avec la projection vue plus haut: 
\begin{center}
   \textit{
      Chercher la distance entre un point \(P\) et une droite contenant \(A\) revient à\+ projeter \(\vv{AP}\) sur un vecteur orthogonal à la droite (en utilisant l'angle complémentaire).
   }
\end{center}
Par ailleurs, de la formule \(\star\) et de la définition du produit vectoriel, on peut en déduire une autre expression de la distance:
\[
   \frac{\vectNorm{\vv{PS} \times \vv{v}}}{\vectNorm{\vv{v}}}   
\]
Sachant l'aire du parallélogramme et un de ses cotés, cette formule consiste exactement à retrouver la hauteur du parallèlogramme, qui constitura ici la distance recherchée. \<

Nous pouvons aussi être amenés à calculer la distance entre un point \(P\) et un plan \(\mathscr{P}\)\+
Si on connait un point \(A\) et un vecteur normal \(\vv{v}\) du plan, alors on peut trouver l'angle entre \(\vv{AP}\) et \(\vv{v}\) et il suffit de calculer la composante scalaire du vecteur \(\vv{AP}\) par rapport au vecteur normal.\+
Visullement cela donne:
\begin{center}
   \begin{tikzpicture}[
      >=stealth,scale=1,line cap=round,
      bullet/.style={circle,inner sep=0.5pt,fill}
   ]
      \coordinate (P1) at (0, 0, 0);
      \coordinate (P2) at (12, 0, 0);
      \coordinate (P3) at (12, 0, -8);  
      \coordinate (P4) at (0, 0, -8);      
    
      \coordinate (X) at (3, 1, -4);

      \coordinate (V1) at (8, 0, -4);
      \coordinate (V2) at (8, 2, -4);

      \coordinate (Proj1) at (8, 1, -4);
      \coordinate (Proj2) at (3, 0, -4);

      \draw[->] (V1) node[below right]{\(A\)} -- (V2) node[above left]{\(\vv{v}\)};
      \draw[->] (V1) -- (X) node[above left]{\(P\)};

      \draw[solid, line width=1pt, color=BrightBlue1] (V1) -- (Proj1);
      \draw node[color=BrightBlue1] at (9.25, 0.5, -4) {\(\vectNorm{\vv{AP}}\cos(\alpha)\)};
      \draw[dashed, color=gray] (X) -- (Proj1);
      \draw[solid, line width=1pt, color=BrightBlue1] (X) -- (Proj2);
      \draw[dashed, color=gray] (Proj2) -- (V1);

      %4,58
      \draw[->] (V1) -- (X) node[above left]{\(P\)};

      \draw[dashed, color=gray] (P1) -- (P2) -- (P3) -- (P4) -- (P1);
      \draw[dashed, color=black!10, fill=BrightBlue1, opacity=0.15] (P1) -- (P2) -- (P3) -- (P4) -- cycle;
      \RightAngle{(10, 0, -4)}{(V1)}{(V2)};
      \RightAngle{(X)}{(Proj2)}{(V1)};

      \draw pic["$\alpha$", draw, gray, angle radius=0.8cm] {angle = V2--V1--X};
      \tkzMarkSegment[color=black, pos=0.5, size=3pt, mark=||](Proj2, X)
      \tkzMarkSegment[color=black, pos=0.5, size=3pt, mark=||](V1, Proj1)

   \end{tikzpicture}   
\end{center}\<

Finalement, il peut être utile de savoir calculer \textbf{l'angle entre deux plans}, il suffit alors de calculer \textbf{l'angle entre leurs vecteurs normaux}.

\chapter*{\chapterstyle{Géométrie {:} Arcs paramétrés}}
Dans le chapitre précédent, en considérant l'espace euclidien \(\R^2\), on a vu qu'on pouvait définir une droite comme étant une \textbf{fonction vectorielle} qui a chaque réel \(t\) associe un point de \(\R^2\) sur la droite.\<

Soit \(I \subset \R\) un intervalle et \(x, y \in \mathscr{F}(I, \R)\), alors on définit la fonction vectorielle suivante:
\[
   \begin{aligned}
      f: I &\longrightarrow \R^2\\
      t &\longmapsto (x(t), y(t))
   \end{aligned}
\]
On apelle alors \textbf{arc paramétré} le couple \((I, f)\) et \textbf{trajectoire ou support} l'ensemble des points \(f(I)\) des points géométriques.\+
Inversement, si on a une courbe \(C \subset \R^2\) et qu'il existe une fonction \(f\) telle que \(C = f(I)\), alors \((I, f)\) est apellé \textbf{paramétrage de la courbe}.\<

Par ailleurs, on dit qu'une fonction vectorielle \(f\) est \textbf{dérivable sur I} (par extension de classe \(\mathcal{C}^k\)) si chacune de ses composantes sont dérivables sur \(I\) (ou de classe \(\mathcal{C}^k\) sur \(I\)) et on note alors \(f'\) sa \textbf{dérivée} qui s'obtient simplement en dérivant composante par composante.\<

Géométriquement \(f'(t)\) représente \textbf{le vecteur tangent} à l'arc parametré au temps \(t\).

\subsection*{\subsecstyle{Propriétés{:}}}
Si \(I\) est un intervalle fermé de la forme \(\icc{a}{b}\), alors on appelle \textbf{point initial} (respectivement \textbf{point final}) le point \(f(a)\) (respectivement \(f(b)\)).\<

Pour un support donné, il existe une infinité d'arc paramétrés décrivant ce support, par exemple les deux arcs paramétrés suivants sont bien différents mais ont \textbf{même support} (le cercle unité).
\begin{align*}
   \lambda_1 &:= \bigl(\icc{0}{2\pi}, (\cos(t), \sin(t))\bigl) \\
   \lambda_2 &:= \bigl(\icc{2\pi}{4\pi}, (\cos(t), \sin(t))\bigl)
\end{align*}
On peut aussi remarquer que tout graphe d'une fonction \(g\) peut être considéré comme le support d'un arc paramétré \((I, f)\) où \(f(t) = (t, g(t))\). \<

Si \(f\) est injective, on dit alors que l'arc paramétré est \textbf{simple}\footnote[1]{Qui consiste à dire que la courbe n'a pas de point multiples}. \+
Si \(f'\) ne s'annule pas, on dit alors que l'arc paramétré est \textbf{régulier}\footnote[2]{Qui consiste à dire que la courbe n'a pas de point stationnaires}. \+

On appelle \textbf{courbe fermée} la trajectoire d'une arc paramétré par une fonction périodique.

\subsection*{\subsecstyle{Courbes polaires et paramétrage par l'angle polaire{:}}}
On fixe une origine \(O\) du plan et une demi-droite partant de l'origine.
Tout point du plan peut alors être identifié par ses \textbf{coordonées polaires} \((r, \theta)\) avec \(r\) la distance dirigée\footnote[3]{ On autorise alors une rayon négatif qui signifie simplement que le point a pour coordonées \((r, \pi + \theta)\), ie le point se situe toujours sur la droite d'angle \(\theta\) mais "de l'autre coté" de l'orgine.} entre \(P\) et l'origine et \(\theta\) l'angle entre \(\vv{OP}\) et la demi-droite.\<

Une façon fréquente de définir une courbe géométrique est alors de donner son \textbf{equation polaire} de la forme \(r = f(\theta)\) pour \(\theta \in I\) qui est alors un cas particulier d'arc paramétré de la forme \((I, (f(\theta)\cos(\theta), f(\theta)\sin(\theta)))\) par passage en coordonées cartésiennes.
\pagebreak
\subsection*{\subsecstyle{Courbes remarquables{:}}}
On peut considèrer alors quelques courbes remarquables:
\begin{figure}[h]
   \centering
   \begin{subfigure}{.3\textwidth}
      \centering
      \begin{tikzpicture}[line cap=round]
         \draw[-latex] (0,0) -- (3.5,0) node [right] {$x$};
         \draw[-latex] (0, 0)     -- (0,1) node [above] {$y$};
         \draw[BrightBlue1, thick] plot[variable=\t,domain=0:4*pi,smooth,thick] ({(\t - sin(\t r))/4},{(1 - cos(\t r))/4});
      \end{tikzpicture}
      \caption*{La cycloïde}
   \end{subfigure}\quad
   \begin{subfigure}{.3\textwidth}
      \centering
      \begin{tikzpicture}[line cap=round]
         \draw[-latex] (-1.2,0) -- (1.2,0) node [right] {$x$};
         \draw[-latex] (0,-1)     -- (0,1) node [above] {$y$};
         \draw[BrightBlue1, thick] plot[variable=\t,domain=0:2*pi,smooth,thick] ({(cos(\t r)+1)*(cos(\t r)/2},{(cos(\t r)+1)*(sin(\t r)/2});
      \end{tikzpicture}
      \caption*{Une cardioïde}
   \end{subfigure}\quad
   \begin{subfigure}{.3\textwidth}
      \centering
      \begin{tikzpicture}[line cap=round]
         \draw[-latex] (-1.2,0) -- (1.2,0) node [right] {$x$};
         \draw[-latex] (0,-1.2)     -- (0,1.2) node [above] {$y$};
         \draw[BrightBlue1, thick] plot[variable=\t,domain=0:2*pi,smooth,thick] ({(cos(\t r)^3)},{(sin(\t r)^3)});
      \end{tikzpicture}
      \caption*{Une hypocycloïde}
   \end{subfigure}\quad
   \begin{subfigure}{.3\textwidth}
      \centering
      \begin{tikzpicture}[line cap=round]
         \draw[-latex] (-0.2,0) -- (1.2,0) node [right] {$x$};
         \draw[-latex] (0,-0.75)     -- (0,0.75) node [above] {$y$};
         \draw[BrightBlue1, thick] plot[variable=\t,domain=0:2*pi,smooth,thick] ({(1/1.8+cos(\t r))/2*cos(\t r)},{(1/2+cos(\t r))/1.8*sin(\t r)});
      \end{tikzpicture}
      \caption*{Un limaçon}
   \end{subfigure}\quad
   \begin{subfigure}{.3\textwidth}
      \centering
      \begin{tikzpicture}[line cap=round]
         \draw[-latex] (-1.2,0) -- (1.2,0) node [right] {$x$};
         \draw[-latex] (0,-0.75)     -- (0,0.75) node [above] {$y$};
         \draw[BrightBlue1, thick] plot[variable=\t,domain=0:2*pi,smooth,thick] ({(sqrt(2)/1.5*sin(\t r))/(1+cos(\t r)^2)},{(sqrt(2)/1.5*sin(\t r)*cos(\t r))/(1+cos(\t r)^2)});
      \end{tikzpicture}
      \caption*{Le lemniscate de Bernoulli}
   \end{subfigure}\quad
   \begin{subfigure}{.3\textwidth}
      \centering
      \begin{tikzpicture}[line cap=round]
         \draw[-latex] (-1.2,0) -- (1.2,0) node [right] {$x$};
         \draw[-latex] (0,-0.75)     -- (0,0.75) node [above] {$y$};
         \draw[BrightBlue1, thick] plot[variable=\t,domain=0:2*pi,smooth,thick] ({\t/7*cos(\t r)},{\t/7*sin(\t r)});
      \end{tikzpicture}
      \caption*{La spirale d'Archimède}
   \end{subfigure}\quad
\end{figure}

Ces courbes ont des propriétés trés intéressantes et se retrouvent, par exemple, en physique. En effet, la cycloïde (inversée) représente la pente qui minimise le temps de trajet d'une bille lancée en haut de celle ci. On dit que c'est une courbe \textbf{brachistochrone}.

\subsection*{\subsecstyle{Longueur de courbes paramétrées{:}}}
Soit \((I, f)\) un arc paramétré \textbf{régulier} de classe \(\mathcal{C}^1\). On s'intéresse à la \textbf{longueur de l'arc paramétré} qu'on notera \(L\).\<

On peut alors subdiviser \(I\) en différents intervalles et approximer la longueur de l'arc par la somme des cordes d'extérmités les différentes subdivisions. En faisant alors tendre la longueur des subdivisions vers 0, on obtient, à la limite, une intégrale de Riemann\footnote[1]{L'intégrande est bien intégrable (et même continue) en tant que somme produit et composées de fonctions intégrables (continues)}, et la définition suivante:
\[
   \mathbox{L := \int_I \vectNorm{f'(t)}\, dt}  
\]
Par exemple dans \(\R^2\) pour l'arc paramétré \((I, f)\) avec \(I = \icc{0}{1}\) et \(f(t) = (t, t^2)\), on obtient:
\[
   L := \int_{I} \vectNorm{f'(t)}\,dt = \int_{\icc{0}{1}} \vectNorm{(1, 2t)}\,dt = \int_{\icc{0}{1}}\sqrt{1 + 4t^2}\,dt \approx 1.47
\]
Quand le support de l'arc paramétré est une partie de \(\R^2\), on peut aussi préfèrer exprimer l'arc en coordonées polaires (ie, \(f(\theta)=(r(\theta)\cos(\theta),r(\theta)\sin(\theta))\)), alors on peut déduire de la définition ci-dessus la caractérisation suivante:
\[
   \mathbox{L := \int_I \sqrt{r'(\theta)^2 + r(\theta)^2}\, d\theta}
\]

\pagebreak
\subsection*{\subsecstyle{Aires de courbes paramétrées{:}}}
Soit \((I, f)\) un arc paramétré de classe \(\mathcal{C}^1\) qui a pour support une partie de \(\R^2\) et \(I\) un segment. On s'intéresse à \textbf{l'aire entre l'arc paramétré} et l'axe des abcisses qu'on notera \(A\).\<

On peut alors subdiviser \(I\) en différents intervalles et approximer l'aire sous la courbe par la somme d'aires de rectangles de largeur \(x(t_i)-x(t_{i-1})\) et de hauteur \(y(x(\overline{t}_i))\) avec \(\overline{t}_i\) une valeur quelconque dans la i-ème subdivision. En faisant alors tendre la longueur des subdivisions vers 0, on obtient une intégrale de Riemann\footnote[1]{\label{note1}L'intégrande est bien intégrable (et même continue) en tant que somme produit et composées de fonctions intégrables (continues)}, et la définition suivante:
\[
   \mathbox{A := \int_I y(t)x'(t)\, dt}  
\]
Par exemple pour l'arc paramétré \((I, f)\) avec \(I = \icc{0}{1}\) et \(f(t) = (t, t^2)\), on obtient:
\[
   A := \int_I y(t)x'(t)\, dt = \int_{\icc{0}{1}} t^2\, dt = \frac{1}{3}
\]

Dans le cas des courbes en coordonées polaires, on peut calculer l'aire bornée par la courbe, on subdivise alors \(I\) en \((\theta_i)_{i \in \N}\), et on calcule la somme de \textbf{secteurs circulaires} de rayon \(r(\theta_k)\). Or on sait que l'aire d'une cercle de rayon \(r\) est donnée par \(\pi r^2\) donc l'aire d'un secteur circulaire d'angle \(\theta\) est donné par \(A = \frac{1}{2}\theta r^2\).\<

Pour obtenir l'aire délimitée par la courbe polaire on peut alors sommer les aires des secteurs circulaires, on obtient alors, à la limite, une intégrale de Riemann et la définition suivante:
\[
   \mathbox{A := \frac{1}{2} \int_I r(\theta)^2\, d\theta}  
\]
\begin{center}
   \textit{
      On a alors deux types de calculs d'aires, l'aire des rectangles entre une courbe et l'axe des abcisses et l'aire des secteurs circulaires entre une courbe et l'origne, qui ne sont \textbf{pas du tout équivalentes}.
   }
\end{center}

\subsection*{\subsecstyle{Reparamétrage{:}}}
Soit \(\gamma = (I, f)\) un arc paramétré de classe \(\mathcal{C}^k\). Soit \(\phi\) un \(\mathcal{C}^k\)-diféomorphisme d'un certain intervalle \(J\) dans \(I\), alors on appelle la fonction \(f \circ \phi\) un \textbf{paramétrage admissible} de l'arc et on dit alors que les arcs sont \(\mathcal{C}^k\)\textbf{-équivalents}.\<

Deux arcs \(\mathcal{C}^k\)\textbf{-équivalents} ont même support, mêmes tangentes, même longueur d'arcs (entre deux points d'une part, et deux image de l'autres) et même points réguliers\footnote[2]{Ceci justifie de contraindre \(\phi\) à être un difféomorphisme.}.


\subsection*{\subsecstyle{Abcisse curviligne et paramétrage normal{:}}}
Soit \(\gamma = (I, f)\) un arc paramétré de classe \(\mathcal{C}^1\). On appelle \textbf{abcisse curviligne} d'origine \(t_0\) la fonction\footnote[3]{Voir 1.}:
\[
   \mathbox{s(t) = \int_{t_0}^{t} \vectNorm{v'(x)} \d x}  
\]
En particulier, on voit que l'abcisse curviligne évaluée en en réel \(t\), est exactement \textbf{la longueur de l'arc} entre \(t_0\) et \(t\), ce qui explique notamment la dénomination anglophone \textit{arc-length function}.\<

Si \(s\) est l'abcisse curviligne d'un arc régulier \(\gamma = (I, f)\) de classe \(\mathcal{C}^1\), alors on peut montrer que \(s\) est un \(\mathcal{C}^1\)-difféomorphisme de \(I\) sur \(J\), et en particulier, on peut alors paramétrer l'arc par l'abcisse curviligne en:
\[
   \mathbox{\gamma' = (J, f \circ s^{-1})}   
\]
On peut alors remarquer alors que ce nouvel arc paramétré s'affranchit de la vitesse de parcours, et parcours la courbe \textbf{à vitesse constante}. Ce reparamétrage appelée \textbf{paramétrage normal} est généralement bien plus naturel quand on s'intéresse à des questions géométriques intrinsèques à la courbe.